% Appendix D — Impossibility & Limits (Phase-5)
% Schema–corollary organization with (i) admissible runs as executions in \NltimesE
% under Phase-0 constraints and (ii) explicit-identity options (syntactic vs normal-form).
% Depends on Appendix A (static axioms) and Appendices B–C (dynamics + collectives).

%\subsection{Impossibility \& Limits}

\subsubsection{Purpose and Scope}
This section collects structural \emph{no-go theorems} that delimit what can be guaranteed within FRFP-like architectures under Explicit--Tacit Separation (ETS), Human-Exclusive Grounding (HEG), and Human-Exclusive Closure (HEC). The results are invariant to model internals, training data, or intent.

The central organizing idea is that many desirable guarantees are \emph{tacit-dependent}: whether they hold can vary with human tacit state even when explicit artifacts are held fixed. Under ETS and AE (AI-Explicit Restriction), tacit-dependent guarantees cannot be enforced or decided by explicit-only mechanisms.

\subsubsection{Standing Assumptions and Notation}

Throughout this appendix we fix the FRFP kernel and assume that all interaction and
epistemic constraints from the main text are in force.

\paragraph{Foundational layers.}
On the single-agent side we assume:
\begin{itemize}
  \item the interaction-layer constraints of Section~2 and Appendix~A.2.6 (admissible traces, no AI
        access to tacit inputs, separation of content-producing and correctness-bearing moves);
  \item the epistemic ontology and explicit/tacit separation from Section~3 and Appendix~A.2--A.3
        (ambient category $C_{\mathrm{full}}$, explicit and tacit subcategories $E$ and $T$, boundary
        functor $\RB : E \to T$, and the six primitives $RI, EC, ED, RB, TE, HFD$);
  \item the explicit workflow algebra and correctness layer from Sections~4--6 and
        Appendix~A.4--A.8 (the TDG-generated explicit category, navigation category $N$, Grothendieck
        semidirect product $N \ltimes E$, semantic functor $\sem{-} : N \ltimes E \to T$, and the
        correctness quotient $\gamma : \Obj(T) \to J$).
\end{itemize}

\paragraph{Dynamic and collective layers.}
In addition, we assume the epistemic dynamics and collective semantics developed in the later parts of
Appendix~A:
\begin{itemize}
  \item the tacit degradation operator $\delta : T \to T$, requirement map
        $\Req : E \times \mathbb{R}_{\ge 0} \to T$, and the induced notion of hallucination and
        finite safe horizons from Appendix~A.9;
  \item the population and consensus layer of Appendix~A.10, which models a population of agents
        with a shared explicit category, local tacit spaces, and a communication graph, and equips
        this structure with a Markovian collective dynamics.
\end{itemize}
For a fixed explicit pipeline $P$ and an admissible implementation of $P$ in this dynamic model,
Appendix~A.9 associates to each finite horizon $N$ a hallucination probability $p_N$ and a
hazard profile $(h_n)_{n \geq 0}$, and defines a safe-horizon functional
$H_P(\varepsilon)$ (see Definition~A.\ref{def:safe-horizon}) capturing the maximal time $N$ for which
the implementation keeps hallucination risk below a tolerance $\varepsilon$. On the population side,
Appendix~A.10 provides analogous notions of collective hallucination time and collective
safe horizons.

The explicit-only impossibility and institutional non-automation results in this appendix are
understood relative to this entire stack: they assume the interaction, epistemic, algebraic,
dynamic, and collective layers just described, and ask which kinds of detectors, consensus
mechanisms, and governance operators can be realised \emph{purely explicitly} within that fixed
architecture.


\subsubsection{Admissible Runs as Interaction Protocol Filtered Executions in \texorpdfstring{N$\ltimes$E}{N ltimes E}}
\begin{definition}[Executions in N$\ltimes$E]
Fix a valid pipeline term $P$ (Epistemic Space / Algebra well-typed). An \emph{execution trace} of $P$ is a finite or infinite sequence
\[
\rho:\quad x_0 \to x_1 \to x_2 \to \cdots
\]
where each $x_k$ is a configuration in N$\ltimes$E reachable from an initial configuration for $P$ by the combined reduction relation $\to$.
\end{definition}

\begin{definition}[Interaction Protocol admissibility predicate]
Let $\mathsf{Adm}_0(\rho)$ be a predicate on execution traces encoding Interaction Protocol constraints:
\begin{enumerate}
  \item \textbf{IL (Interaction Limitation):} $\rho$ is turn-based and respects externally imposed bounds (e.g. finite-length turns).
  \item \textbf{AR (Alignment Restriction):} explicit actions occur only in response to explicit instructions present in the trace.
  \item \textbf{MD (Mode Discipline):} no step in $\rho$ applies tacit operators inside the explicit phase; tacit transitions occur only after the boundary.
  \item \textbf{NTER/CSC:} no step reconstructs tacit evaluation from explicit traces or treats explicit computation as a correctness oracle.
\end{enumerate}
We treat $\mathsf{Adm}_0$ as an axiomatically specified filter on traces.
\end{definition}

\begin{definition}[Admissible runs]
An \emph{admissible run} of a pipeline $P$ is an execution trace $\rho$ in N$\ltimes$E such that $\mathsf{Adm}_0(\rho)$ holds.
We write $\mathrm{Runs}(P)$ for the set of admissible runs of $P$.
\end{definition}

\begin{remark} (Admissible runs and explicit-only results).
This definition pins all explicit-only impossibility and non-automation statements to a common
operational semantics: admissible runs are exactly those executions in \(N \ltimes E\) that
additionally satisfy the interaction constraints (IL, AR, MD, NTER/CSC). All explicit-only
no-go theorems in Appendix~A.11 are understood to quantify over this class of runs.

\end{remark}

\subsubsection{Two Notions of ``Same Explicit Artifact''}
No-go statements often quantify over pairs of runs that produce the ``same explicit artifact.'' There are two natural notions.

\begin{definition}[Syntactic explicit identity]
Two runs are \emph{syntactically explicit-identical} if they produce equal explicit artifacts as elements of the explicit object $E$ (i.e. equality in the representation domain).
\end{definition}

\begin{definition}[Semantic explicit identity via normal forms]
Assume the explicit reduction system $\to_{\Ecat}$ is terminating and locally confluent (Appendix~A axioms), so each explicit artifact has a unique normal form $\mathrm{nf}(e)$.
Two runs are \emph{normal-form explicit-identical} if their produced artifacts $e,e'$ satisfy
\[
\mathrm{nf}(e)=\mathrm{nf}(e').
\]
\end{definition}

\begin{remark}
Syntactic identity is appropriate when artifacts are literal strings/terms. Normal-form identity is appropriate when explicit equivalence is taken modulo rewriting (e.g. definitional equalities, normalization, canonical forms). All no-go results below hold under either identity notion; when needed we indicate which is assumed.
\end{remark}
\begin{remark}[Admissible runs and stochastic trajectories]
In the dynamic semantics of Appendix~B, the same pipeline $P$ gives rise
to a measurable space of trajectories $\Omega_{\mathrm{allowed}}$,
together with a family of probability measures capturing different
stochastic implementations. We implicitly identify each admissible run
$\rho$ of $P$ in the combined system $\mathsf{N} \ltimes \mathsf{E}$
with its corresponding trajectory
$\omega \in \Omega_{\mathrm{allowed}}$. In particular, the hallucination time $\tau(\omega)$, the finite-horizon probabilities $(p_N)$,
the hazard sequence $(h_n)$, and the safe-horizon function $H_P(\varepsilon)$
of Appendix~B are all computed over this common set of admissible runs.

\end{remark}

\subsubsection{Tacit Dependence via \texorpdfstring{$\gamma(\sem{P})$}{gamma(sem(P))}}
\begin{definition}[Tacit-dependent predicates]
Let $P$ be a fixed valid pipeline and let $\mathrm{Runs}(P)$ be its admissible runs.
A predicate $\mathcal{G}$ about an explicit artifact is \emph{tacit-dependent} if there exist two admissible runs $\rho_1,\rho_2\in\mathrm{Runs}(P)$ such that:
\begin{itemize}
  \item the runs produce the same explicit artifact (either syntactically or up to normal form, per Section~D.3), and
  \item the observable correctness outcomes differ:
  \[
    \gamma(\sem{P})\ \text{differs between}\ \rho_1\ \text{and}\ \rho_2.
  \]
\end{itemize}
\end{definition}

\begin{remark}
This definition formalizes ``depends on tacit context'' using only Epistemic Algebra semantics: $\sem{-}$ captures tacit transformation induced by the pipeline run, and $\gamma$ extracts the human-observable correctness class.
\end{remark}
\begin{remark}[Examples of tacit-dependent predicates]
\label{rem:tacit-dependent-examples}
The definition above applies not only to pointwise correctness
judgements but also to quantitative long-horizon properties. Typical
examples of tacit-dependent predicates $G$ include:
\begin{itemize}
  \item the event that a run hallucinates within a fixed horizon $N$,
        i.e.\ $G(\omega)$ is the predicate $[\tau(\omega) \le N]$;
  \item the statement that an implementation has $\varepsilon$-safe horizon at least $N$, i.e.
  $G$ is the predicate $[H_P(\varepsilon) \ge N]$ for some fixed $\varepsilon \in (0,1)$;

  \item in the collective setting of Appendix~A.10, the statement that the
        collective hallucination time exceeds a fixed horizon, i.e.\
        $G$ is the predicate $[\tau_{\mathrm{coll}} > N]$ for a given
        population and consensus dynamics.
\end{itemize}
In each case, there can exist admissible runs with the same explicit
artifacts (in the sense of Definition~B.\,66) but different values of
$\tau$, $\HP(\varepsilon)$ or $\tau_{\mathrm{coll}}$, so $G$ is
tacit-dependent in the sense of Definition~B.\,66.
\end{remark}

\subsubsection{Schema Theorem: Tacit-Dependent Guarantees Are Not Explicit-Implementable}
\begin{definition}[Explicit-only mechanism]
An \emph{explicit-only mechanism} is any procedure definable using explicit morphisms in $\Ecat$ (and finite control) whose inputs and outputs live entirely in explicit space. Concretely, one may model an explicit-only mechanism as a total function $M:E^k\to E^\ell$, or as a family of such functions closed under composition.
\end{definition}

\begin{theorem}[Tacit-dependent guarantees are not explicit-implementable]\label{thm:D-schema}
Assume ETS and AE. Let $\mathcal{G}$ be a tacit-dependent predicate about explicit artifacts (Section~D.4), where ``same explicit artifact'' is interpreted either syntactically or by normal form (Section~D.3). Then no explicit-only mechanism can enforce $\mathcal{G}$ for all admissible runs.
\end{theorem}

\begin{proof}[Proof sketch]
By tacit dependence, there exist admissible runs producing the same explicit artifact (under the chosen identity notion) but yielding different outcomes $\gamma(\sem{P})$. Any explicit-only mechanism sees only the explicit artifact(s) and therefore cannot act differently on these runs. Hence it cannot enforce a guarantee whose truth varies across admissible runs while explicit data are held fixed.
\end{proof}

\subsubsection{Corollaries (No-Go Family)}
Each of the following is an instantiation of Theorem~\ref{thm:D-schema}.

\subsubsubsection{No Perfect Explicit-Only Hallucination Detector}
\begin{definition}[Explicit-only monitor]
An \emph{explicit-only monitor} is a function $M:E\to\{\mathrm{ok},\mathrm{flag}\}$.
\end{definition}

\begin{definition}[Perfect detector]
$M$ is \emph{perfect} if, for every admissible run producing explicit artifact $e$ with tacit context $t$ at evaluation time,
\[
M(e)=\mathrm{ok}\Rightarrow t\not\prec\mathsf{Req}(e),\qquad
M(e)=\mathrm{flag}\Rightarrow t\prec\mathsf{Req}(e).
\]
\end{definition}

\begin{corollary}[No perfect explicit-only detector]\label{cor:D-no-perfect-detector}
Assume ETS and AE, and assume requirements are nontrivial and contexts vary across admissible runs. Then no explicit-only monitor is a perfect hallucination detector.
\end{corollary}

\begin{proof}[Proof sketch]
The predicate ``$e$ is a hallucination'' is tacit-dependent: the same explicit artifact (syntactically or up to normal form) can be within grounding in one admissible run and beyond grounding in another. Apply Theorem~\ref{thm:D-schema}.
\end{proof}

\begin{corollary}[No elimination of hallucination under degradation]
\label{cor:no-elimination-under-degradation}
Assume the Epistemic Dynamics conditions that yield almost-sure
hallucination inevitability: in particular, a finite-capacity tacit
state space, a monotone non-invertible degradation operator
$\delta \colon T \to T$, and recurring nontrivial requirements as in
Theorem~B.\,29. Let $\tau$ denote the hallucination time random variable
for a fixed pipeline $P$ and any admissible stochastic implementation
of $P$ whose trajectories lie in $\Omega_{\mathrm{allowed}}$.
Then there is no explicit-only policy which guarantees
\[
  \tau(\omega) \;=\; \infty
\]
for all admissible runs $\omega \in \Omega_{\mathrm{allowed}}$, nor even
$\mathbb{P}(\tau = \infty) = 1$ for all such implementations.
\end{corollary}

\begin{proof}
Under the assumptions of Appendix~B, Theorem~B.\,29 states that every
admissible stochastic implementation of $P$ has hallucination time
$\tau$ finite almost surely:
\[
  \mathbb{P}(\tau < \infty) \;=\; 1
\]
whenever its support is contained in $\Omega_{\mathrm{allowed}}$.
Consider the predicate
\[
  G(\omega) \;\equiv\; [\,\tau(\omega) = \infty\,]
\]
on admissible runs. By construction of the dynamic semantics, there
exist admissible runs with the same explicit artifacts but different
tacit evolutions along which $\tau$ takes different values; in
particular, $G$ is tacit-dependent in the sense of
Definition~B.\,66. If there were an explicit-only policy that guaranteed
$\tau = \infty$ almost surely for all admissible implementations, that
policy would enforce the tacit-dependent predicate $G$ by purely
explicit means, contradicting Theorem~B.\,69. Hence no explicit-only
policy can eliminate hallucination under the Epistemic Dynamics degradation
assumptions.
\end{proof}
\begin{corollary}[No explicit-only infinite safe horizon]\label{cor:no-infinite-safe-horizon}
Assume the Epistemic Dynamics conditions of Appendix~B that yield almost-sure hallucination
inevitability, and let $H_P(\varepsilon)$ be the $\varepsilon$-safe horizon of a
stochastic FRFP pipeline $P$ as in Definition~B.41. Then for any fixed
$\varepsilon \in (0,1)$ there is no explicit-only policy that guarantees
$H_P(\varepsilon) = \infty$ for all admissible stochastic implementations of $P$.
\end{corollary}

\begin{proof}
Under the Epistemic Dynamics assumptions, Theorem~B.29 implies that for any admissible
implementation we have $P(\tau < \infty) = 1$, so for every $\varepsilon \in (0,1)$
the set $\{N \in \mathbb{N} : p_N \le \varepsilon\}$ is bounded almost surely, i.e.\
$H_P(\varepsilon) < \infty$ along admissible runs. The predicate
$[H_P(\varepsilon) = \infty]$ is tacit-dependent (Remark~B.69), so an explicit-only
policy enforcing it for all admissible implementations would contradict
Theorem~B.69.
\end{proof}


\subsubsubsection{No Fully Automated Correctness Oracle}
\begin{definition}[Correctness oracle]
A \emph{correctness oracle} is a function $\Omega:E\to J$ (e.g. $J=\{\mathrm{accept},\mathrm{reject},\mathrm{revise}\}$) such that for every admissible run of a pipeline $P$ producing artifact $e$,
\[
\Omega(e)\;=\;\gamma(\sem{P}).
\]
\end{definition}

\begin{corollary}[No explicit correctness oracle]\label{cor:D-no-oracle}
Assume HEG and HEC. No explicit-only oracle $\Omega$ can reproduce $\gamma(\sem{P})$ for all admissible runs without violating the FRFP axioms.
\end{corollary}

\begin{proof}[Proof sketch]
Correctness outcomes $\gamma(\sem{P})$ are tacit-dependent by HEG/HEC: they are determined by tacit evaluation and closure, which can vary across admissible runs even when the explicit artifact is held fixed (under the chosen explicit identity). Apply Theorem~\ref{thm:D-schema}.
\end{proof}

\subsubsubsection{No Grounding Gain from Explicit Consensus}
Assume the functorial consensus model of Appendix~A.10.

\begin{corollary}[Consensus cannot create grounding]\label{cor:D-no-consensus-grounding}
Let $\mathcal{C}:\Ecat^n\to\Ecat$ be any consensus functor acting purely in explicit space. For any agent $i$ with tacit state $t_i$, if each input artifact $e_k$ exceeds grounding for $i$, then $\mathcal{C}(e_1,\dots,e_n)$ also exceeds grounding for $i$.
\end{corollary}

\begin{proof}[Proof sketch]
The predicate ``within grounding for agent $i$'' is tacit-dependent. Since $\mathcal{C}$ is explicit-only, it cannot enforce the guarantee that aggregation restores grounding. Apply Theorem~\ref{thm:D-schema}.
\end{proof}

\subsubsubsection{No Fully Automated Consensus Stability Guarantee}
\begin{corollary}[No automated consensus guarantee]\label{cor:D-no-auto-consensus}
Under ETS and HEG/HEC, no fully automated consensus pipeline can guarantee that its outputs are free of hallucination across all admissible runs and populations.
\end{corollary}
In the Markovian population semantics of Appendix~A.10, any choice of
consensus functors and explicit-only coordination rules determines a
collective dynamics $M$ on the product category together with a
collective hallucination time $\tau_{\mathrm{coll}}$. Under the Epistemic Dynamics
assumptions lifted to the population level, Theorem~B.\,58 shows that
$\tau_{\mathrm{coll}}$ is almost surely finite and, moreover, that its
almost-sure finiteness is invariant under insertion of additional
explicit-only consensus layers. A purported fully automated consensus
pipeline guaranteeing hallucination-free outputs for all admissible
populations would therefore enforce the tacit-dependent predicate
$[\tau_{\mathrm{coll}} = \infty]$ using only explicit mechanisms,
contradicting Theorem~B.\,69 together with Theorem~B.\,58.

\begin{proof}[Proof sketch]
A guarantee of hallucination-freedom is tacit-dependent and would yield a correctness oracle or perfect detector for consensus outputs. Apply Theorem~\ref{thm:D-schema} and Corollary~\ref{cor:D-no-perfect-detector}.
\end{proof}

\subsubsubsection{Catching Does Not Imply Eliminating}
\begin{corollary}[Detection does not imply elimination]\label{cor:D-detect-vs-elim}
Even if a subclass of failures is detectable by explicit invariants (e.g. syntactic checks, type checks, cross-consistency tests), this does not imply that hallucination (defined by tacit adequacy) is eliminable in general.
\end{corollary}

\begin{proof}[Proof sketch]
Explicit detectability concerns explicit predicates on $E$. Hallucination is tacit-dependent (Section~D.6.1). Eliminating hallucination would require enforcing a tacit-dependent guarantee by explicit-only means, contradicting Theorem~\ref{thm:D-schema}.
\end{proof}

\subsubsection{Interpretation and Design Consequences}
The no-go theorems do not claim that automation is useless; they claim that automation cannot replace tacit grounding without changing axioms.

Mitigations must therefore be expressed as:
\begin{itemize}
  \item restricting what is demanded of explicit space (reducing $\mathsf{Req}$),
  \item injecting fresh tacit evaluation (additional humans, replication),
  \item designing institutions that slow confidence inflation and feedback (Appendix~E).
\end{itemize}
Appendix~B refines these impossibility results with quantitative
structure. Each stochastic implementation of a pipeline induces a
hallucination time $\tau$, a finite-horizon risk curve
$(p_N)_{N \in \mathbb{N}}$, a hazard sequence $(h_n)_{n \in \mathbb{N}}$
and, when defined, a safe-horizon functional $\HP(\varepsilon)$. The
dynamic refinement order of Definition~B.\,43 then compares pipeline
implementations by their finite-horizon risks. Within the structural
limits above, design work can still aim to shape hazard profiles so that
risk is concentrated outside the operational horizon of interest and to
select implementations that dynamically refine a baseline, even though
eventual hallucination and the tacit-dependence of correctness remain
unavoidable.

Appendix~A.10 extends the picture to populations and collective dynamics.
A choice of consensus operators and coordination policies determines a
Markov process $M$ on the space of explicit population states and a
collective hallucination time $\tau_{\mathrm{coll}}$. The collective
inevitability and invariance results of Appendix~10  show
that purely explicit-only changes to consensus and logging cannot alter
the inevitability class of $\tau_{\mathrm{coll}}$. What they can do is
change the collective hazard profile and the distribution of
$\tau_{\mathrm{coll}}$, for example by producing a dynamics $M_2$ that
is a population-level dynamic refinement of some baseline $M_1$. In
this sense, long-horizon design remains meaningful: it can bend but not
break the structural bounds established in this appendix.

% End Appendix D
